\sgaStableTarget{sga6.e08.prop.3-2}\textbf{Proposition 3.2.} \uline{Let} @@M0@@ \uline{be a projective morphism}, \uline{a complete intersection}, \uline{whose virtual relative dimension} @@M1@@ \uline{is constant}. \uline{Suppose that} @@M2@@ \uline{is quasi-compact}, \uline{and has an ample sheaf}. \uline{Then for every} @@M3@@ \uline{we have}:
@@M4@@

We thus obtain, in particular, an additive homomorphism:
@@M5@@
sending @@M6@@ into @@M7@@.

\sgaStableTarget{sga6.e08.para.3-3}\textbf{3.3.} Recall that in (V 6.1) we defined a functor @@M8@@ on the category of graded @@M9@@-algebras, with values in the category of augmented @@M10@@-@@M11@@-algebras, where @@M12@@ is a binomial ring (V 2.7). If @@M13@@ is a graded @@M14@@-algebra, we have:
@@M15@@

For every augmented @@M16@@-@@M17@@-algebra @@M18@@, we defined a @@M19@@-@@M20@@-homomorphism, called the ``completed Chern class'' (V 6.7):
@@M21@@

Finally, we defined a ring homomorphism, functorial in @@M22@@, called the Chern character (V 6.4):
@@M23@@
and a group endomorphism, called the Todd character (V 1.16):
@@M24@@
